(1) 两片接触的薄透镜在空气中，参考波长处光焦度分别为 $\Phi_1=1/f_1$、$\Phi_2=1/f_2$（可正可负），整体光焦度 $\Phi=\Phi_1+\Phi_2$。给定各片材料的Abbe数 $V_1,V_2>0$ 且 $V_1\ne V_2$。说明近轴薄透镜接触时光焦度为何相加。
(2) 采用一级色差模型，两个规定谱线之间各片光焦度变化满足 $\Delta\Phi_i\simeq\Phi_i/V_i$，推导色差抵消条件。
(3) 在指定总体光焦度 $\Phi$ 下解出 $\Phi_1,\Phi_2$ 的解析表达式。
(4) 若 $V_1>V_2$、$\Phi>0$，判断两片透镜的正负焦度并解释实际剩余二级色谱及厚透镜间隔效应。
